\begin{figure}[!htbp]
\centering
\resizebox{\textwidth}{!}{%
\begin{tikzpicture}[
    >=Stealth,
    box/.style={draw, rounded corners=2pt, minimum height=0.7cm, minimum width=1.4cm,
                font=\small, align=center, thick},
    databox/.style={box, fill=blue!8},
    procbox/.style={box, fill=orange!12},
    sortbox/.style={box, fill=green!12},
    votebox/.style={box, fill=red!10},
    arr/.style={->, thick},
    lbl/.style={font=\footnotesize, text=black!70},
]
% === Shared encoding steps (top row); every view expands and standardizes the same way ===
\node[databox] (input) at (0,0) {$x \in \R^d$};
\node[procbox, right=0.6 of input] (poly) {Polynomial\\[-1pt]Expansion};
\node[procbox, right=0.6 of poly] (std) {Standardization};
\draw[arr] (input) -- (poly);
\draw[arr] (poly) -- (std);
\node[lbl, above=0.15 of input] {$d$ feat.};
\node[lbl, above=0.15 of poly] {$d'$ feat.};
% === One projection and argsort per view: view k ranks the standardized features with its own W_k ===
\node[procbox] (proj1) at (0,-1.75) {Projection $W_1$,\\[-1pt]$\argsort$};
\node[procbox] (proj2) at (4,-1.75) {Projection $W_2$,\\[-1pt]$\argsort$};
\node[procbox] (projK) at (8,-1.75) {Projection $W_K$,\\[-1pt]$\argsort$};
\node at ($(proj2)!0.5!(projK)$) {\Large$\cdots$};
\draw[arr] (std.south) -- ++(0,-0.45) -| (proj1.north);
\draw[arr] (std.south) -- ++(0,-0.45) -| (proj2.north);
\draw[arr] (std.south) -- ++(0,-0.45) -| (projK.north);
% === Multi-view branching (3 networks), each reading its own input ranking ===
\node[sortbox, minimum height=1.1cm] (net1) at (0, -3.6) {ArrowFlow\\[-1pt]Net 1};
\node[sortbox, minimum height=1.1cm] (net2) at (4, -3.6) {ArrowFlow\\[-1pt]Net 2};
\node[sortbox, minimum height=1.1cm] (netK) at (8, -3.6) {ArrowFlow\\[-1pt]Net $K$};
\node at ($(net2)!0.5!(netK)$) {\Large$\cdots$};
\draw[arr] (proj1) -- node[lbl, right] {$\pi_1 \in S_e$} (net1);
\draw[arr] (proj2) -- node[lbl, right] {$\pi_2 \in S_e$} (net2);
\draw[arr] (projK) -- node[lbl, right] {$\pi_K \in S_e$} (netK);
% === Predictions ===
\node[databox, below=0.6 of net1] (pred1) {$\hat{y}_1$};
\node[databox, below=0.6 of net2] (pred2) {$\hat{y}_2$};
\node[databox, below=0.6 of netK] (predK) {$\hat{y}_K$};
\draw[arr] (net1) -- (pred1);
\draw[arr] (net2) -- (pred2);
\draw[arr] (netK) -- (predK);
% === Majority vote ===
\node[votebox, below=0.7 of pred2, minimum width=2.8cm] (vote) {Majority vote};
\draw[arr] (pred1.south) -- ++(0,-0.2) -| (vote.160);
\draw[arr] (pred2) -- (vote);
\draw[arr] (predK.south) -- ++(0,-0.2) -| (vote.20);
\node[databox, below=0.5 of vote, fill=red!15] (final) {$\hat{y}$};
\draw[arr] (vote) -- (final);
% === Network detail inset (far right, well separated) ===
\begin{scope}[shift={(13.8,-1.9)}]
\node[font=\small\bfseries, anchor=south] at (0,0.5) {Network Detail};
\draw[thick, rounded corners=3pt, fill=green!4] (-3.35,-4.85) rectangle (2.85,0);
\coordinate (insetleft) at (-3.35,-2.4);
\node[sortbox, minimum width=2.6cm, font=\footnotesize] (L1) at (0,-0.55) {Ranking layer 1\\[-1pt]$N_1$ filters};
\node[sortbox, minimum width=2.6cm, font=\footnotesize] (L2) at (0,-1.95) {Ranking layer 2\\[-1pt]$N_2$ filters};
\draw[arr] (L1) -- node[lbl, right, font=\scriptsize] {$\pi'{\in}S_{N_1}$} (L2);
\node[lbl, above=0.08 of L1, font=\scriptsize] {$\pi_k \in S_e$};
% below the last hidden layer: the output layer, used in training, and the nearest-neighbor readout, used to predict
\node[sortbox, minimum width=2.1cm, fill=green!20, font=\scriptsize] (Lout) at (-1.25,-3.55) {Output layer $L$\\[-1pt]$C$ filters};
\node[votebox, minimum width=2.1cm, font=\scriptsize] (RO) at (1.25,-3.55) {kNN readout\\[-1pt]stored rankings};
\draw[arr, dashed] ([xshift=-4mm]L2.south) -- ++(0,-0.4) -| (Lout.north);
\draw[arr] ([xshift=4mm]L2.south) -- ++(0,-0.4) -| (RO.north);
% training: the motions returned by the output votes train the last hidden layer (dashed, upward), which relays signed
% motions to the layer below (dashed, between the two hidden layers)
\draw[dashed, thick] (Lout.west) -- (-2.85,-3.55) -- (-2.85,-1.95);
\draw[arr, dashed] (-2.85,-1.95) -- (L2.west);
\node[lbl, font=\scriptsize, rotate=90, anchor=south] at (-2.85,-2.75) {returned motions};
\draw[arr, dashed] ([xshift=-7mm]L2.north) -- node[lbl, left, font=\scriptsize] {relay} ([xshift=-7mm]L1.south);
\node[lbl, font=\scriptsize, anchor=east] at (-1.33,-2.3) {$\pi''{\in}S_{N_2}$};
\node[lbl, below=0.08 of Lout, font=\scriptsize, align=center] {training:\\[-1pt]votes};
\node[lbl, below=0.08 of RO, font=\scriptsize, align=center] {prediction: vote of the\\[-1pt]nearest stored rankings};
\end{scope}
% Dashed connector from network box to inset
\draw[dashed, gray, thick] (netK.east) -- ++(0.8,0) |- (insetleft);
\end{tikzpicture}%
}
\caption{\textbf{ArrowFlow pipeline.} A real-valued input $x$ is encoded by polynomial expansion of degree $p$, standardization, a projection and argsort. The $K$ views use different projections $W_k$, $k = 1, \dots, K$, cycling through target-aware, random and calibrated strategies, so their networks receive different input rankings. Each network (detail, right) is a stack of ranking layers, and every layer ranks its ranking filters by Spearman's footrule distance to its input ranking. In training, an output layer holds one filter per class, and the filter of the example's class votes to move toward the last hidden ranking. The motions this vote returns train the last hidden layer, which relays signed motions to the preceding layer (dashed arrows). To predict, a nearest-neighbor (kNN) readout compares the last hidden ranking with the stored hidden rankings of the training examples. The $K$ predictions are combined by majority vote. $S_n$ denotes the set of rankings of $n$ items.}
\label{fig:pipeline}
\end{figure}
